%% SODE-Guard — USENIX Security '27 submission (problem-space evaluation).
%% Anonymous for double-blind review. Conference ID: SEC27.
%%
%% This is the EXPERIMENTAL (problem-space) companion to the feature-space work.
%% It is a DIFFERENT contribution: it measures whether the robustness rankings
%% hold when evasions are carried out as real network traffic. It cites the
%% detector in the third person and uses no results from the other paper.
%%
%% Drop in the official template before submission:
%%   https://secartifacts.github.io / https://www.usenix.org/conferences/author-resources
%% replace the preamble below with \documentclass[...]{usenixsecurity}.
%% Results tables marked %% FROM TESTBED are filled by testbed/analysis/make_tables.py
%% once the Hetzner runs complete; no number is written here by hand.
\documentclass[letterpaper,twocolumn,10pt]{article}
\usepackage{usenix-preamble}   % local shim; swap for the official class
\usepackage{graphicx}
\usepackage{booktabs}
\usepackage{xcolor}
\usepackage{hyperref}
\usepackage[capitalise,noabbrev]{cleveref}
\newcommand{\todo}[1]{\textcolor{red}{[\textbf{TODO:} #1]}}
\newcommand{\fromtestbed}[1]{\textcolor{blue}{#1}}  %% replaced by real numbers

\begin{document}

\title{\Large \bf Does Robustness Survive the Wire? A Problem-Space Evaluation of
Certified Neural Intrusion Detectors}
\author{}         % anonymous
\date{}
\maketitle
\thispagestyle{empty}

\begin{abstract}
Adversarially trained and certified neural network intrusion detectors are
reported to resist evasion in \emph{feature space}: an attacker who perturbs the
extracted flow statistics within a budget cannot flip the prediction. It is far
less clear that these guarantees mean anything on the wire, where an attacker
controls packets, not features, and every realisable change moves many features
at once while having to keep the attack working. We build an isolated testbed in
which standard evasions---padding, timing, packet splitting, chaff and
rate shaping---are carried out as real traffic against deliberately vulnerable
containerised services, and we re-extract features and re-score the resulting
flows. A functionality oracle discards any evasion that breaks the attack, so
every reported evasion is a working one. On this problem-space benchmark we
compare adversarially trained neural detectors, a randomized-smoothing
certificate, gradient-boosted trees, a deployed detection platform, and the
signature engines Suricata and Zeek, under a gradient-guided and a black-box
problem-space attack. We find that \fromtestbed{the feature-space rankings
\{persist/change\}}: \todo{headline result, from testbed}. We release the
testbed, the attack orchestration, and every captured flow.
\end{abstract}

\section{Introduction}
Learned intrusion detectors are evaluated almost entirely in feature space.
Robustness papers, including the certified ones, perturb the vector of flow
statistics and report that the prediction is stable. But a network attacker does
not edit a feature vector; they send packets, and a tool like CICFlowMeter then
summarises those packets into features. Two gaps follow. First, not every
feature is under the attacker's control, and the ones that are move together when
a real knob (padding, delay, splitting) is turned. Second, an evasion is only
interesting if the attack still works; a perturbation that silences the exploit
is a failed attack, not a successful evasion. Neither gap is visible in a
feature-space experiment.

This paper asks a single question: \emph{do the robustness rankings established
in feature space survive when the evasion is realised as traffic?} We answer it
with a measurement study rather than a new model. Our contributions:
\begin{itemize}
\item \textbf{A problem-space testbed} (\cref{sec:testbed}) that runs benign load
and seven attack families against vulnerable containers on an isolated cluster,
captures traffic, extracts the same features used to train the detectors, and
labels flows by \mbox{(attacker, window)}.
\item \textbf{Two problem-space attacks} (\cref{sec:attack}): a gradient-guided
attack that maps the constrained feature-space adversarial direction onto
physical traffic knobs, and a budgeted black-box search that needs only detector
scores, so it runs identically against neural detectors, trees, the platform and
the signature engines.
\item \textbf{A functionality oracle} (\cref{sec:oracle}) that certifies each
evasion still achieves its goal, making evasion-success rates meaningful.
\item \textbf{A fair comparison} (\cref{sec:results}) of adversarially trained
neural detectors, a smoothing certificate, trees, a deployed platform, and
Suricata/Zeek on identical realised traffic, with a measured deployment cost.
\item \textbf{An open release} (\cref{sec:openscience}) of the testbed, attack
code and captured flows.
\end{itemize}

\section{Background and threat model}\label{sec:threat}
\paragraph{Detectors under test.} We evaluate a neural stochastic differential
equation detector with adversarial training and a randomized-smoothing
certificate, described in prior work~\cite{anon-featurespace} and used here
unmodified as one of several detectors; matched adversarially trained MLP and
neural-ODE baselines; gradient-boosted trees (LightGBM, XGBoost, random forest);
a deployed detection platform exposing a scoring API; and the signature engines
Suricata (ET Open) and Zeek. We treat all of them as black boxes that map a flow
to a score, except where an attack is explicitly white-box.

\paragraph{Attacker.} The attacker controls a host inside the monitored network
and wants its attack traffic classified as benign while the attack still works.
It may change traffic only through realisable knobs (\cref{sec:attack}), each with
a budget, and must keep the attack functional (\cref{sec:oracle}). It may be
white-box (gradient-guided) or black-box (search over knobs using scores only).
The defender observes flows and raises an alarm at a chosen false-alarm budget.

\section{The problem-space testbed}\label{sec:testbed}
\cref{fig:testbed} shows the testbed. An isolated k3d cluster hosts a
microservice shop, Node-RED, Grafana and InfluxDB (versions with public CVEs),
with Locust generating benign user sessions throughout. An attacker pod holds
standard public tools. We capture one pcap per run on the cluster bridge, extract
flows with a single pinned CICFlowMeter build (the same extractor family as the
training corpus), and label each flow from the attacking host and the run window.
\emph{Isolation.} The cluster runs on an internal-only network; a host firewall
rule drops any forwarded traffic toward the public interface, a per-namespace
network policy denies external egress, and a target guard aborts any scenario
whose target is not a private address. Full detail in \cref{sec:ethics}.

\begin{figure}[t]\centering
\fbox{\parbox{0.9\columnwidth}{\centering\small victims + benign load + attacker
$\rightarrow$ pcap $\rightarrow$ CICFlowMeter $\rightarrow$ flows $\rightarrow$
\{neural, trees, platform, Suricata/Zeek\}}}
\caption{Testbed data flow. \todo{replace with the real architecture figure.}}
\label{fig:testbed}
\end{figure}

\paragraph{Scenarios.} Seven families mirror the training corpus and the common
public datasets: reconnaissance (nmap), vulnerability scan (nuclei), SSH/HTTP
brute force (hydra), slow and volumetric DoS (slowloris, rate-limited hping3),
and three exploits (Grafana path traversal CVE-2021-43798, InfluxDB auth bypass
CVE-2019-20933, Node-RED RCE). Benign traffic runs alongside every scenario.

\paragraph{Two tracks.} \emph{Track A} (primary) trains every detector on a clean
testbed corpus and attacks the trained detectors with fresh evasive runs, so
train and test share one environment and domain shift cannot explain the result.
\emph{Track B} scores the detectors trained on the external container corpus on
the testbed traffic without retraining, measuring cross-environment transfer.

\section{Problem-space attacks}\label{sec:attack}
\paragraph{Transformations.} Each evasion is a vector of budgeted knobs
(\cref{tab:knobs}): payload padding ($\le{+}50\%$ bytes), timing delay and jitter
($\le 50$\,ms/packet, $\le 3\times$ duration), packet splitting via MSS clamping
($\le 4\times$ packets), chaff ($\le 1$ benign request per attack request) and
rate shaping. Timing and rate are enforced in the kernel with \texttt{tc};
padding, splitting and chaff are honoured by the HTTP scenarios. Budgets are
hard-capped so no search can exceed them.

\begin{table}[t]\centering\small
\begin{tabular}{lll}
\toprule
knob & mechanism & budget\\
\midrule
padding & junk header bytes & $\le{+}50\%$ bytes\\
timing & \texttt{tc}/netem delay+jitter & $\le 50$\,ms/pkt, $\le 3\times$\\
splitting & MSS clamp & $\le 4\times$ packets\\
chaff & interleaved benign requests & $\le 1{:}1$\\
rate & egress rate cap & attack stays functional\\
\bottomrule
\end{tabular}
\caption{Problem-space transformation budgets.}\label{tab:knobs}
\end{table}

\paragraph{Gradient-guided.} We compute the constrained feature-space adversarial
direction on the target (white box) over the mutable features, then map the sign
of each feature's gradient to the knob that can physically move it (e.g.\ forward
packet length $\to$ padding, inter-arrival time $\to$ delay). We realise the knob,
re-extract and re-score.

\paragraph{Black-box.} A budgeted $(1{+}\lambda)$ search over the knob vector
minimises the target's margin on the realised flows, using scores only, subject
to the budget and the oracle ($\approx 50$ queries per scenario and detector).
Because it needs only scores, the identical search runs against the trees, the
platform and the signature engines.

\section{Functionality oracle}\label{sec:oracle}
An evasive run counts only if the attack still works: brute force still finds a
valid credential; an exploit still reads the file or writes the RCE marker; a scan
still reports the same open ports; a DoS still raises benign p95 latency by at
least $2\times$. The oracle reads the artefacts each scenario leaves and gates
every reported evasion, so evasion-success rates never include broken attacks.

\section{Evaluation}\label{sec:results}
\paragraph{Protocol.} Five repetitions $\times$ five seeds. Detection is reported
per flow and per attack instance (a run is detected if any of its flows alarms),
with macro-F1 on realised attacks. Alarms use conformal thresholds calibrated on
benign testbed traffic at 1\% and 5\% false-alarm budgets; we report the test
false-alarm rate alongside detection. Contrasts are declared in advance and
tested with a paired bootstrap under Holm correction:
(i)~neural adversarially trained vs.\ MLP adversarially trained,
(ii)~neural vs.\ trees, (iii)~neural vs.\ Suricata, each under both attacks.

\paragraph{Q1: do the rankings persist?} \todo{Table from testbed: detection and
macro-F1 under realised attacks, per detector.} %% FROM TESTBED detection.md
\begin{table}[t]\centering\small
\begin{tabular}{lcccc}
\toprule
detector & det.\ rate & macro-F1 & FAR@1\% & FAR@5\%\\
\midrule
neural (AT) & \fromtestbed{--} & \fromtestbed{--} & \fromtestbed{--} & \fromtestbed{--}\\
MLP (AT) & \fromtestbed{--} & \fromtestbed{--} & \fromtestbed{--} & \fromtestbed{--}\\
trees & \fromtestbed{--} & \fromtestbed{--} & \fromtestbed{--} & \fromtestbed{--}\\
platform & \fromtestbed{--} & \fromtestbed{--} & \fromtestbed{--} & \fromtestbed{--}\\
Suricata & \fromtestbed{--} & \fromtestbed{--} & \fromtestbed{--} & \fromtestbed{--}\\
\bottomrule
\end{tabular}
\caption{Detection on realised evasive traffic (Track A).}\label{tab:detection}
\end{table}

\paragraph{Q2: can the attacker evade while the attack works?} \todo{Evasion
success and overhead, gradient vs.\ black-box.} %% FROM TESTBED evasion.md

\paragraph{Q3: deployment cost.} \todo{Per-flow latency, throughput, memory on the
server CPU for every detector, Suricata and the platform.} %% FROM TESTBED cost.md

\paragraph{Q4: signatures and the platform.} \todo{Suricata/Zeek and platform
alerts on the same pcaps vs.\ the neural detectors.}

\paragraph{Track B (transfer).} \todo{Detectors trained on the container corpus,
scored on testbed traffic without retraining.}

\section{Related work}\label{sec:related}
Problem-space adversarial examples and the feature/problem-space gap
\cite{pierazzi2020problemspace}; realisable network evasions and feature
mutability \cite{apruzzese2023realworld}; adversarial training and randomized
smoothing \cite{madry2018,cohen2019}; evaluation pitfalls in ML security
\cite{arp2022dodo,carlini2019evaluating}. \todo{position against the detector's
own feature-space paper and recent problem-space IDS work.}

\section{Ethics and compliance}\label{sec:ethics}
This is defensive measurement. All attacks are standard and public and run only
against our own containers on an isolated VM, separate from any production host.
No attack traffic can leave the host: the cluster is on an internal-only network,
a host firewall rule drops forwarded traffic to the public interface, a network
policy denies external egress, and a target guard aborts any non-private target.
Floods are rate-limited. The data are flow statistics only, with no payloads or
personal identifiers. We disclosed nothing to third parties because no external
system is involved. Following USENIX guidance, we used large-language-model
assistance for copy-editing and for parts of the released code; the authors
checked every derivation, ran every experiment, and verified every citation.

\section{Open science}\label{sec:openscience}
We release the testbed manifests, the attack and evasion orchestration, the
functionality oracle, the scoring and signature harness, the analysis code, and
every captured flow (features only) and run manifest, under an open licence at an
anonymous repository for review (\url{https://anonymous.4open.science/r/SEC27-sode-guard}).
All tables are regenerated from the released runs by a single script.

{\footnotesize \bibliographystyle{plain} \bibliography{sode_guard_usenix}}

\end{document}
